\documentclass[hyperref={unicode=true}]{beamer}
\usepackage[T2A]{fontenc}
\usepackage[utf8]{inputenx}
\usepackage[russian]{babel}

\usepackage{multicol}

\usepackage[pgf]{dot2texi}
\usepackage{tikz}
\usetikzlibrary{shapes, arrows}

\usepackage{listings}
\usepackage{graphicx}

\usepackage{comment}

\title{Лекция 6. Среда Anis для динамической верификации на соответствие формальной модели Event-B}
\author{}
\date{}

\usetheme{Warsaw}

\AtBeginSection[] {
	\begin{frame}{Содержание}
		\tableofcontents[currentsection]
	\end{frame}
}
%\overfullrule=5pt

\begin{document}
	\begin{frame}{}
		\titlepage
	\end{frame}

    \begin{frame}{Цель лекции}
    Познакомиться с тем, как можно проводить динамическую верификацию
    ядра Linux на соответствие формальной модели и какую поддержку
    для этого оказывает среда Anis.
    \end{frame}

    \section{Введение}

	\begin{frame}{Для чего нужен Anis}
    \begin{itemize}
    \item Позволяет задействовать формальную модель
            Event-B как оракул для динамической верификации
            на соответствие этой модели.
    \item Возможно и иное использование оракула.
    \item На данный момент динамическая верификация выполняется
            на сценариях, описываемых тестами.
    \item В планах применить этот оракул для реального
            мониторинга ядра.
	\end{itemize}
    \end{frame}

    \begin{frame}{Основной компонент Anis}
    \begin{itemize}
    \item это транслятор Event-B в Python
    \item Позволяет задействовать формальную модель
            как оракул в тестах на Python
    \item pytest -- инфраструктура тестирования на Python.
    \end{itemize}
    \end{frame}

    \begin{frame}{Как запустить Anis}
    \begin{itemize}
    \item Это расширение для \texttt{VSCode}
    \item Дополняет \texttt{VSCode} инструментами для
        разработки моделей Event-B и их использования
    \end{itemize}
    \end{frame}

    \section{Тестирование при помощи pytest}

    \begin{frame}{Определение теста в pytest}
    \begin{itemize}
    \item Тест --- это функция или метод класса с пустым конструктором
    \item Тест пройден, если функция завершается нормально
    \item Тест провален, если функция выбрасывает исключение \texttt{AssertionError}
    \item Тест ошибочен, если функция выбрасывает другое исключение
    \item Поэтому проверки в тесте ожидаются при помощи оператора \texttt{assert}
    \end{itemize}
    \end{frame}

    \begin{frame}{Упражнение}
    В модуле \texttt{int2text.py} написана ошибочная реализация функции,
    которая разбивает число от 0 до 999 на сумму чисел (возвращается как кортеж).
    Если прочитать числа в этом кортеже последовательно, должно получиться
    исходное число.

    Напишите тесты для этой функции и найдите все ошибки в ее реализации.
    \end{frame}

    \begin{frame}{Что еще есть в pytest}
    \begin{itemize}
    \item параметризованные тесты (помечены декоратором pytest.mark.parametrize)
    \item фиксчи, в том числе session-scoped
    \item фиксчи с финализацией при помощи yield
    \item параметризованные фиксчи
    \item плагины с предопределенными фиксчами
    \item \texttt{conftest.py} для фиксч, видимых для всех тестов
    \end{itemize}
    \end{frame}

    \begin{frame}{VSCode}
    \begin{itemize}
    \item \texttt{VSCode} --- расширяемая IDE (частично open-source)
    \item Расширения для Python
    \item Команды
    \item Настраивание расширений и среды (в том числе, на уровне User и Workspace)
    \item Multi-root workspace
    \item Запуск и отладка, их конфигурирование
    \item Task'и
    \item Активация расширений
    \end{itemize}
    \end{frame}

    \begin{frame}{Тестирование в VSCode}
    \begin{itemize}
    \item Необходима виртуальная среда python
    \item Вкладка Testing: дерево тестов, запуск тестов и отладка
    \item Тестирование со сбором покрытия по коду
    \item Расширение Coverage Gutters: показ покрытия в редакторе
    \item Генерация HTML отчета о покрытии кода (\texttt{--cov-report html})
    \end{itemize}
    \end{frame}

    \section{Тестирование ядра Linux}

    \begin{frame}{Объект тестирования}
    \begin{itemize}
    \item Цель --- тщательно протестировать модуль безопасности
    \item Тестирование модульное? системное? системное нацеленное!
    \item Доступ до модуля безопасности осуществляется в
            системных вызовах (проверка возможности предоставить доступ)
    \end{itemize}
    \end{frame}

    \begin{frame}{Методика тестирования}
    \begin{itemize}
    \item Документ, где требования к модулю безопасности изложены
        в виде тестовых ситуаций
    \item Содержит предположения и ограничения, в которых проводится
        тестирование
    \end{itemize}
    \end{frame}

    \begin{frame}{Структура теста}
    \begin{itemize}
    \item setup --- подготовка системы (файлы, пользователи, конфигурация ядра)
    \item run --- осуществление целевого системного вызова (или последовательности
                    системных вызовов) некоторым исполнимым модулем
    \item teardown --- восстановление системы (необязательно, если тест выполняется
            в контейнере и не задействует конфигурацию ядра хоста)
    \end{itemize}
    \end{frame}

    \begin{frame}{LinuxSpec}
    \begin{itemize}
    \item Такая структура требует разных инструментов: скрипт для записи setup,
            Си для записи программы с целевыми системными вызовами, Python
            для запуска тестов
    \item Разным тестам может потребоваться разные способы их выполнения
    \item Поэтому удобно найти способ записи, который бы скрыл эти подробности
    \item \texttt{LinuxSpec} --- интерфейс компонента, который выполняет шаги
        на этапе setup, превращает описание цепочки системных вызовов в Си
        программу и ее компилирует, может по-разному выполнять тест
    \end{itemize}
    \end{frame}

    \begin{frame}{Как выглядит LinuxSpec}
    \begin{itemize}
    \item конкретные действия зависят от методики тестирования. Например (\texttt{t: LinuxSpec}):
    \item \texttt{t.make\_file(path, owner, group, mode)} --- создать файл
    \item \texttt{with t.make\_prog() as child: child.open(path, flags)}
        --- превратить описание цепочки системных вызовов (описывается
                как последовательность методов объекта \texttt{child})
            в исполнимый модуль
    \item \texttt{t.run(prog, user, group, umask)} --- запустить программу
        под мониторингом и собрать трассу для перепроверки
    \item Надо описывать LinuxSpec вручную для каждого проекта 
    \end{itemize}
    \end{frame}

    \begin{frame}{Упражнение}
    Ознакомьтесь с примером проекта по тестированию ядра
    согласно методике тестирования. Найдите отсутствующий тест,
    требуемый по методике, и реализуйте его.
    \end{frame}

    \section{Anis}

    \begin{frame}{Общие сведения}
    \begin{itemize}
    \item Анис -- это инструмент для работы с моделями Event-B
	    на основе их трансляции в исполнимую программу (Python, C, ...)
    \item Цель -- скорость воспроизведения трассы
    \item Как следствие, Анис не использует солверы (будет медленное воспроизведение)
    \item Доступен как расширение VSCode
    \item Активация: по открытию файлов \texttt{.bm} (машина),
	    \texttt{.bc} (контекст), вызову команды \texttt{"Event-B: ..."}
		или при активации расширения Python
    \item Работа в multi-root workspace не рекомендуется (плохо работает Pylance)
    \item Директория src в workspace -- для модели
    \end{itemize}
    \end{frame}

    \begin{frame}{Работа с моделью}
    \begin{itemize}
    \item В Анис есть текстовый редактор для машин и контекстов
    \item Автодополнение, раскраска синтаксиса, outline, сниппеты
    \item Импорт/экспорт модели со средой Rodin (можно вести разработку
          параллельно в двух средах)
    \item Синтаксический анализатор, проверка статической типизации 
    \item Анис транслирует только одну из машин (\texttt{componentUnderTest}), ее нужно выбрать
    \item Способы выбора componentUnderTest: Settings, командой,
          codeLens "Set as componentUnderTest" прямо в редакторе
    \end{itemize}
    \end{frame}

    \begin{frame}{Представление модели в Python}
    \begin{itemize}
    \item Модель транслируется в директорию model в workspace
    \item Плоская модель -- результат соединения componentUnderTest
          со всеми ее машинами-предками и всеми нужными контекстами
          (получается один компонент со всем)
    \item \texttt{machine.py}: содержит класс \texttt{Machine}
          с плоской моделью, объект представляет модельное состояние
    \item Поля класса \texttt{Machine} -- это константы контекстов
          и переменные машины
    \item Константы (множества) хранятся целиком в памяти, поэтому
          они создаются пустыми и дополняются по мере необходимости в новых элементах
    
    \end{itemize}
    \end{frame}

    \begin{frame}{Представление модели в Python (2)}
    \begin{itemize}
    \item \texttt{test\_axioms.py}: аксиомы превращаются в тестовые
          функции от объекта класса \texttt{Machine}
    \item \texttt{test\_invariants.py}: инварианты аналогично
    \item \texttt{test\_initialisation.py}: действия события
          \texttt{INITIALISATION} аналогично
    \item Остальные события превращаются в отдельные функции, которые
          вычисляют охранные условия
    \end{itemize}
    \end{frame}

    \begin{frame}{Транслируемое подмножество Event-B}
    \begin{itemize}
    \item Без редко используемых операций
    \item Неявные присваивания произвольного вида разрешены
          только в \texttt{INITIALISATION}
    \item В остальных местах они должны иметь вид условного выражения
          \texttt{v :| (pred \& v' = expr1) or (not pred \& v' = expr2)}
    \item Кванторы должны начинаться с явного указания
          области значений подкванторной переменной (для перебора в цикле):
          \texttt{forall x . x in \{....\} => pred}, но не \texttt{forall x . pred}
    \item Итерация по константам-множествам возможна, но только
          по построенной на данный момент части значений 
    \end{itemize}
    \end{frame}

    \begin{frame}{Проверка готовности машины к работе}
    \begin{itemize}
    \item Нужно проверить, что машина правильно оттранслирована
          и проинициализирована
    \item Признак этого: все сгенерированные тесты успешно выполняются
    \item Вкладка Testing, в ней 3 вида тестов (для аксиом, для
          инвариантов и для события \texttt{INITIALISATION})
    \item Если не все тесты проходят, нужно вручную задать правильные
          начальные значения нужных полей
    \item Если тест выполняется долго или не может быть выполнен в принципе
          (т.к. аксиома/инвариант были добавлены в модель для целей
          дедуктивной верификации), то его приходится не запускать
    \end{itemize}
    \end{frame}

    \begin{frame}{Задание начальных значений полей машины}
    \begin{itemize}
    \item Анис не генерирует возможную интерпретацию значений
          констант и начальных значений переменных машины, т.к.
          для этого потребуется солвер
    \item Пользователь должен сам задать эти значения
    \item В файле \texttt{conftest.py} надо найти фиксчу \texttt{m}
          и добавить присваивания во все нужные поля объекта-машины \texttt{m}
    \item Множества должны быть \texttt{frozenset}
    \end{itemize}
    \end{frame}

    \begin{frame}{Медиатор}
    \begin{itemize}
    \item Если уровень абстракции модели выше уровня абстракции в трассе
          мониторинга, требуется компонент-медиатор
    \item Он строит модельную трассу по трассе мониторинга
    \item Для медиатора отводится директория \texttt{mediator}
    \item Основной класс -- \texttt{TraceTranslator}
    \end{itemize}
    \end{frame}

    \begin{frame}{Отладка}
    \begin{itemize}
    \item Расхождения между реализацией и моделью проявляются
          в виде выбрасываемых исключений
    \item Если запустить тест в режиме Debug с точками останова (например,
          на исключениях), то можно остановиться на месте ошибки и
          выполнять отладку привычными средствами
    \item Объект класса \texttt{TraceTranslator} можно
          обернуть при помощи функций сохранения трасс в файл
          (удобно после теста посмотреть отдельно на трассу и даже ее запустить
          без мониторинга --- см. контекстное меню у тестов в Testing)
    \item Для запуска трасс без мониторинга удобно создать launch
    \item В нем нужно прописать \texttt{env: PYTHONPATH}, равную \texttt{\$\{workspaceFolder\}}.
    \end{itemize}
    \end{frame}

    \begin{frame}{Сбор покрытия по модели}
    \begin{itemize}
    \item 5 критериев покрытия:
		\begin{itemize}
                \item строки,
                \item ветви (только охранные условия),
                \item ветви (все элементарные
          условия охранных условий),
                \item mcdc (только охранные условия),
                \item mcdc (все элементарные
          условия охранных условий)
                \end{itemize}
    \item Сбор покрытия нужно включить в Settings
    \item По умолчанию каждый запуск тестов заменяет полностью покрытие
          (можно в Settings это поменять)
    \end{itemize}
    \end{frame}

    \begin{frame}{Просмотр покрытия}
    \begin{itemize}
    \item Просмотр покрытия по модели прямо в редакторе
          при помощи расширения \texttt{Coverage Gutters}
    \item Покрытие строится автоматически после завершения тестов
    \item Есть команда для принудительного построения покрытия без перезапуска тестов
          (необходимо, если поменяли критерий покрытия в Settings)
    \item Есть генерация покрытия с раскраской элементарных условий в виде HTML
    \item Можно создать Task для открытия отчета HTML без повторной генерации
    \end{itemize}
    \end{frame}

    \begin{frame}{Улучшение покрытия}
    \begin{itemize}
    \item Для полного покрытия нужно, чтобы условия реализовались
          в тестах со значениями как \texttt{True}, так и \texttt{False}
    \item Не все охранные условия нужно покрывать целиком:
            \begin{itemize}
            \item Это условие обеспечивается медиатором (например, оно задает тип параметра)
            \item Это условие является теоремой
            \item Тесты не создают ситуации, когда это условие
                  принимает оба значения (например, условие о том,
                  что количество файлов не превышает лимита)
            \end{itemize}
    \item Чтобы показать, что охранное условие не нужно покрывать, нужно
          вначале комментария к этому условию написать дефис (а если в этом
          комментарии есть что-то еще, то после дефиса пробел)
    \end{itemize}
    \end{frame}

\end{document}

